% Generated draft: requires mathematical and source review.
\documentclass[11pt]{article}
\usepackage[margin=25mm]{geometry}
\usepackage{fontspec}
\setmainfont{Latin Modern Roman}
\usepackage{amsmath,amssymb,mathtools,unicode-math}
\setmathfont{Latin Modern Math}
\usepackage{enumitem,needspace,xurl,hyperref}
\setlength{\parindent}{0pt}
\setlength{\parskip}{5pt}
\setlength{\emergencystretch}{3em}
\newcommand{\N}{\mathbb{N}}
\newcommand{\Z}{\mathbb{Z}}
\newcommand{\Q}{\mathbb{Q}}
\newcommand{\R}{\mathbb{R}}
\newcommand{\C}{\mathbb{C}}
\newcommand{\F}{\mathbb{F}}
\newcommand{\A}{\mathbb{A}}
\newcommand{\PP}{\mathbb{P}}
\newcommand{\E}{\mathbb{E}}
\newcommand{\eps}{\varepsilon}
\newcommand{\dd}{\,\mathrm{d}}
\begin{document}
\section*{UnsolvedMath ID: 3800018}
\textbf{Source problem code:} catalogue entry
\subsection{Definitions and mathematical statement}
Fix an integer $d\geq3$. The input is a collection of $n$ axis-aligned boxes
\[
B_i=\prod_{j=1}^d[a_{ij},b_{ij}]\subset\R^d,\qquad a_{ij}\leq b_{ij},\quad 1\leq i\leq n.
\]
Klee's measure problem asks to compute exactly
\[
\operatorname{vol}_d\left(\bigcup_{i=1}^n B_i\right),
\]
where $\operatorname{vol}_d$ is $d$-dimensional Lebesgue measure. The source asks for the optimal worst-case running time as $n\to\infty$, with $d$ fixed. In dimension three it asks in particular whether there is a near-linear algorithm or an $\Omega(n^{3/2})$ lower bound, and asks whether restrictions to fat boxes or equal-sized boxes permit faster algorithms.

A usual fatness condition bounds $\max_j(b_{ij}-a_{ij})/\min_j(b_{ij}-a_{ij})$ by a fixed constant for every positive-volume box. This is an indicated convention only: the source gives no fatness parameter or definition. Likewise, equal-sized could mean translates of one box or another size restriction; equal volume is a distinct condition. The coordinate/computation model, treatment of randomization, and meaning of near-linear are not specified.

The supplied research summary reports $O(n^{3/2})$ in three dimensions and $O(n^{d/2})$ in fixed dimension, and conditional lower bounds of exponent $3/2-o(1)$ under a 3-uniform hyperclique hypothesis, with extensions to higher dimensions. Conditional exponent optimality is not an unconditional $\Omega(n^{3/2})$ lower bound. The original algorithm/lower-bound alternatives and the restricted-box variants must therefore be distinguished from this supplied partial-progress assessment.
\subsection{Short English statement}
Determine the optimal time to compute the volume of a union of axis-aligned boxes in fixed dimension at least three, including faster algorithms for fat or equal-sized boxes. The supplied progress report gives conditional exponent optimality, leaving unconditional and restricted-input questions separate.
\subsection{Sources}
\begin{itemize}
\item Local UnsolvedMath registry, permanent ID 3800018, exact source code catalogue entry, ``Klee's measure problem''; proposed by Jeff Erickson and cited proposers. Registry status partially solved, classification PARTIAL-PROGRESS.
\item Jeff Erickson, \emph{Open Problems in Computational Geometry}, ``Klee's measure problem'', \url{https://jeffe.cs.illinois.edu/open/}. Supplied HTML extraction accessed July 29, 2026. The source warns that its pages have not been systematically updated since 2001.
\item The supplied research summary cites K\"unnemann, FOCS 2022, for a three-dimensional conditional lower bound under the 3-uniform hyperclique hypothesis, and a SoCG 2023 higher-dimensional generalization. No titles, full author lists, theorem hypotheses, or formulation of that hypothesis are provided. The summary says unconditional optimality and fat/equal-box questions are not fully closed, while unit-hypercube cases allow faster algorithms. These supplied claims are not independently verified here.
\item Provenance flags status and rights NEEDS\_REVIEW: public access was verified, later systematic source updates are absent, and redistribution permission was not established.
\end{itemize}
\end{document}
